% addendum-electrical.tex
% Document 2 - Technical Discovery Addendum: Electrical
%
% Build (from this folder, requires XeLaTeX -- see README.md):
%   xelatex -interaction=nonstopmode addendum-electrical.tex
%   xelatex -interaction=nonstopmode addendum-electrical.tex   (rerun for TOC/page refs)
%
% Issued alongside the Pre-Project Intake and Discovery Form (Document 1)
% only when electrical engineering assistance was indicated. See
% docs/Engineering_Project_Intake_Plan.md for the full workflow.

\documentclass[11pt]{article}

\usepackage{../../shared/intake-style}
\usepackage{../../shared/form-components}

\SetDocIdentity{Technical Discovery Addendum -- Electrical}{ENG-FORM-02E}{v0.1-draft}
\SetFieldPrefix{addelec}

% Replace these placeholders once branding and markings are confirmed
% (see "Remaining Organization-Specific Decisions" in the intake plan).
\SetCompanyName{[Engineering Center Name]}
\SetMarkingText{FOR OFFICIAL USE ONLY}

\SetDocPurpose{%
This addendum collects the electrical-engineering detail that the
Pre-Project Intake and Discovery Form (Document 1) does not ask for. Fill
it out only if electrical assistance was indicated on that form. Its
purpose is to give the center's electrical engineers enough detail about
power, controls, and existing electrical systems to scope the work
accurately. It is fine to mark an item Unknown -- the center will help
fill in technical gaps during discussion.%
}

\begin{document}

\begin{Form}

\CoverPage

\DocumentControlBlock

\LinkedIntakeBlock

\IntakeSection{Power and Load Requirements}
\TextArea{PowerLevels}{Voltage levels, phase, and load requirements (known or estimated)}{\linewidth}{1.5cm}\\[8pt]
\TextArea{PowerSource}{Power source and availability (existing service, generator, outage windows)}{\linewidth}{1.5cm}

\IntakeSection{Controls and Instrumentation}
\TextArea{ControlsDescribe}{Controls, instrumentation, or automation involved}{\linewidth}{1.5cm}\\[8pt]
\TextArea{ElecDrawings}{Existing one-line diagrams, panel schedules, or electrical drawings available}{\linewidth}{1.2cm}

\IntakeSection{Hazardous Locations and Classifications}
\TextArea{HazLocation}{Known hazardous-location classifications or explosion-proof requirements}{\linewidth}{1.5cm}

\IntakeSection{Applicable Standards}
\Guidance{List any electrical codes or standards already known to apply
(e.g. NEC, NFPA 70E). Leave Confirmed By / Status blank if not yet
reviewed by a discipline lead -- this records applicability, not
compliance.}\\[4pt]
\StandardsTable{ElecStd}

\IntakeSection{Site and Access Constraints}
\TextArea{SiteAccess}{Site access, outage windows, or lockout/tagout constraints specific to this work}{\linewidth}{1.5cm}

\IntakeSection{Assumptions and Open Questions}
\OpenItemsTable{ElecOpen}

\AcknowledgementBlock{This addendum supplements the Pre-Project Intake and
Discovery Form and reflects current understanding only. It does not commit
\CompanyName's resources and is not authorization to begin engineering
work or to charge project time.}

\end{Form}

\end{document}
